% GENERATED FILE. Edit canonical lessons, then run npm run generate:books.
% canonical-source-hash: fnv1a64:784a9f5914c61c59
% canonical-lessons: ML-C159-katrika, ML-C159-currika, ML-C159-ani, ML-C159-veli, ML-C159-gerru, ML-R159-first-pass-the-house-and-the-town, ML-R159-second-pass-nature-and-food

\chapter{Scissors, Hammer, Nail, Fence, Gate}
\label{ch:ml-scissors-hammer-nail-fence-gate}
\hlchaptermodality{159}

\begin{chapteropening}
\textbf{By the end of this chapter:} \emph{I can say scissors, a hammer, a nail, a fence and a gate.}

Complete the last lesson of chapter 159: I can say scissors, a hammer, a nail, a fence and a gate.
\end{chapteropening}

\section[\texorpdfstring{\ml{കത്രിക} (katrika)}{കത്രിക (katrika)}]{\ml{കത്രിക} (katrika) — scissors}

\label{lesson:ML-C159-katrika}



\begin{quote}
Before the new one: say the Malayalam for a blanket, then the Malayalam for a bucket.
\end{quote}

\subsection*{You'll want to know: \ml{കത്രിക}}
\textbf{\ml{കത്രിക}} — \emph{katrika} — \textquotedblleft{}scissors\textquotedblright{}.

\begin{grammarlens}[title={What you've built}]
One more word for this chapter.
\end{grammarlens}

\subsection*{Your turn}
\begin{itemize}
\raggedright
  \item \emph{Say it:} \emph{katrika}
  \item \emph{Say it:} \emph{katrika}, once more
  \item \emph{Say it:} say \emph{katrika}
  \item \emph{Recall:} say the Malayalam for a blanket, then the Malayalam for a bucket, then say \emph{katrika} again
\end{itemize}

\begin{tcolorbox}[breakable,colback=teal!4,colframe=teal!35!black,title={Before you move on}]
What does \ml{കത്രിക} mean? (\textquotedblleft{}Scissors\textquotedblright{}.) Say it once more.
\end{tcolorbox}
\section[\texorpdfstring{\ml{ചുറ്റിക} (cuṟṟika)}{ചുറ്റിക (cuṟṟika)}]{\ml{ചുറ്റിക} (cuṟṟika) — a hammer}

\label{lesson:ML-C159-currika}



\begin{quote}
Before the new one: say the Malayalam for a bucket, then the Malayalam for scissors.

Which weekday names match Tamil almost exactly? (\emph{Tiṅkaḷ}, \emph{veḷḷi}, \emph{ñāyar}.)
\end{quote}

\subsection*{You'll want to know: \ml{ചുറ്റിക}}
\textbf{\ml{ചുറ്റിക}} — \emph{cuṟṟika} — \textquotedblleft{}a hammer\textquotedblright{}.

\begin{grammarlens}[title={What you've built}]
One more word for this chapter.
\end{grammarlens}

\subsection*{Your turn}
\begin{itemize}
\raggedright
  \item \emph{Say it:} \emph{cuṟṟika}
  \item \emph{Say it:} \emph{cuṟṟika}, once more
  \item \emph{Say it:} say \emph{cuṟṟika}
  \item \emph{Recall:} say the Malayalam for a bucket, then the Malayalam for scissors, then say \emph{cuṟṟika} again
\end{itemize}

\begin{tcolorbox}[breakable,colback=teal!4,colframe=teal!35!black,title={Before you move on}]
What does \ml{ചുറ്റിക} mean? (\textquotedblleft{}A hammer\textquotedblright{}.) Say it once more.
\end{tcolorbox}
\section[\texorpdfstring{\ml{ആണി} (āṇi)}{ആണി (āṇi)}]{\ml{ആണി} (āṇi) — a nail}

\label{lesson:ML-C159-ani}



\begin{quote}
Before the new one: say the Malayalam for scissors, then the Malayalam for a hammer.
\end{quote}

\subsection*{You'll want to know: \ml{ആണി}}
\textbf{\ml{ആണി}} — \emph{āṇi} — \textquotedblleft{}a nail\textquotedblright{}.

\begin{grammarlens}[title={What you've built}]
One more word for this chapter.
\end{grammarlens}

\subsection*{Your turn}
\begin{itemize}
\raggedright
  \item \emph{Say it:} \emph{āṇi}
  \item \emph{Say it:} \emph{āṇi}, once more
  \item \emph{Say it:} say \emph{āṇi}
  \item \emph{Recall:} say the Malayalam for scissors, then the Malayalam for a hammer, then say \emph{āṇi} again
\end{itemize}

\begin{tcolorbox}[breakable,colback=teal!4,colframe=teal!35!black,title={Before you move on}]
What does \ml{ആണി} mean? (\textquotedblleft{}A nail\textquotedblright{}.) Say it once more.
\end{tcolorbox}
\section[\texorpdfstring{\ml{വേലി} (vēli)}{വേലി (vēli)}]{\ml{വേലി} (vēli) — a fence}

\label{lesson:ML-C159-veli}



\begin{quote}
Before the new one: say the Malayalam for a hammer, then the Malayalam for a nail.

Say black, white and blue. (\emph{Karuppŭ}, \emph{veḷḷa}, \emph{nīla}.) Which one is a Sanskrit loan? (\emph{Nīla}.)
\end{quote}

\subsection*{You'll want to know: \ml{വേലി}}
\textbf{\ml{വേലി}} — \emph{vēli} — \textquotedblleft{}a fence\textquotedblright{}.

\begin{grammarlens}[title={What you've built}]
One more word for this chapter.
\end{grammarlens}

\subsection*{Your turn}
\begin{itemize}
\raggedright
  \item \emph{Say it:} \emph{vēli}
  \item \emph{Say it:} \emph{vēli}, once more
  \item \emph{Say it:} say \emph{vēli}
  \item \emph{Recall:} say the Malayalam for a hammer, then the Malayalam for a nail, then say \emph{vēli} again
\end{itemize}

\begin{tcolorbox}[breakable,colback=teal!4,colframe=teal!35!black,title={Before you move on}]
What does \ml{വേലി} mean? (\textquotedblleft{}A fence\textquotedblright{}.) Say it once more.
\end{tcolorbox}
\section[\texorpdfstring{\ml{ഗേറ്റ്} (gēṟṟŭ)}{ഗേറ്റ് (gēṟṟŭ)}]{\ml{ഗേറ്റ്} (gēṟṟŭ) — a gate}

\label{lesson:ML-C159-gerru}



\begin{quote}
Before the new one: say the Malayalam for a nail, then the Malayalam for a fence.
\end{quote}

\subsection*{You'll want to know: \ml{ഗേറ്റ്}}
\textbf{\ml{ഗേറ്റ്}} — \emph{gēṟṟŭ} — \textquotedblleft{}a gate\textquotedblright{}.

\begin{grammarlens}[title={What you've built}]
One more word for this chapter.
\end{grammarlens}

\subsection*{Your turn}
\begin{itemize}
\raggedright
  \item \emph{Say it:} \emph{gēṟṟŭ}
  \item \emph{Say it:} \emph{gēṟṟŭ}, once more
  \item \emph{Say it:} say \emph{gēṟṟŭ}
  \item \emph{Recall:} say the Malayalam for a nail, then the Malayalam for a fence, then say \emph{gēṟṟŭ} again
\end{itemize}

\begin{tcolorbox}[breakable,colback=teal!4,colframe=teal!35!black,title={Before you move on}]
What does \ml{ഗേറ്റ്} mean? (\textquotedblleft{}A gate\textquotedblright{}.) Say it once more.
\end{tcolorbox}
\section[(dialogue)]{The house and the town — a first pass}

\label{lesson:ML-R159-first-pass-the-house-and-the-town}



\begin{quote}
Say the words for a fence and a gate.
\end{quote}

\subsection*{Your turn}
Say these aloud:

\begin{itemize}
\raggedright
  \item a horse, a donkey, a buffalo, a pig, a lion
  \item a rat, a frog, a tortoise, a peacock, a parrot
  \item a dove, a butterfly, an ant, a mosquito, a shoulder
  \item an elbow, a knee, a tongue, skin, a forehead
\end{itemize}

\begin{tcolorbox}[breakable,colback=teal!4,colframe=teal!35!black,title={Before you move on}]
A horse? (\textbf{\ml{കുതിര}}, \emph{kutira}.) A donkey? (\textbf{\ml{കഴുത}}, \emph{kaḻuta}.) A buffalo? (\textbf{\ml{പോത്ത്}}, \emph{pōttŭ}.) A pig? (\textbf{\ml{പന്നി}}, \emph{panni}.) A lion? (\textbf{\ml{സിംഹം}}, \emph{simhaṁ}.) A rat? (\textbf{\ml{എലി}}, \emph{eli}.) A frog? (\textbf{\ml{തവള}}, \emph{tavaḷa}.) A tortoise? (\textbf{\ml{ആമ}}, \emph{āma}.) A peacock? (\textbf{\ml{മയിൽ}}, \emph{mayil}.) A parrot? (\textbf{\ml{തത്ത}}, \emph{tatta}.) A dove? (\textbf{\ml{പ്രാവ്}}, \emph{prāvŭ}.) A butterfly? (\textbf{\ml{പൂമ്പാറ്റ}}, \emph{pūmpāṟṟa}.) An ant? (\textbf{\ml{ഉറുമ്പ്}}, \emph{uṟumpŭ}.) A mosquito? (\textbf{\ml{കൊതുക്}}, \emph{kotukŭ}.) A shoulder? (\textbf{\ml{തോൾ}}, \emph{tōḷ}.) An elbow? (\textbf{\ml{കൈമുട്ട്}}, \emph{kaimuṭṭŭ}.) A knee? (\textbf{\ml{കാൽമുട്ട്}}, \emph{kālmuṭṭŭ}.) A tongue? (\textbf{\ml{നാവ്}}, \emph{nāvŭ}.) Skin? (\textbf{\ml{തൊലി}}, \emph{tolī}.) A forehead? (\textbf{\ml{നെറ്റി}}, \emph{neṟṟi}.)
\end{tcolorbox}
\section[(dialogue)]{Nature and food — a second pass}

\label{lesson:ML-R159-second-pass-nature-and-food}



\begin{quote}
Say the words for a cheek and a chin.
\end{quote}

\subsection*{Your turn}
Say these aloud:

\begin{itemize}
\raggedright
  \item a cheek, a chin, a street, a library, a shirt
  \item trousers, an earring, a bangle, a ring, a cup
  \item a pot, a stove, a hand fan, an electric fan, a bicycle
  \item a boat, a train, a pillow, a blanket, a bucket
  \item scissors, a hammer, a nail, a fence, a gate
\end{itemize}

\begin{tcolorbox}[breakable,colback=teal!4,colframe=teal!35!black,title={Before you move on}]
A cheek? (\textbf{\ml{കവിൾ}}, \emph{kavil}.) A chin? (\textbf{\ml{താടി}}, \emph{tāṭi}.) A street? (\textbf{\ml{തെരുവ്}}, \emph{teruvŭ}.) A library? (\textbf{\ml{ഗ്രന്ഥശാല}}, \emph{granthaśāla}.) A shirt? (\textbf{\ml{ഷർട്ട്}}, \emph{ṣarṭṭŭ}.) Trousers? (\textbf{\ml{പാന്റ്}}, \emph{pāṇṟŭ}.) An earring? (\textbf{\ml{കമ്മൽ}}, \emph{kammal}.) A bangle? (\textbf{\ml{വള}}, \emph{vaḷa}.) A ring? (\textbf{\ml{മോതിരം}}, \emph{mōtiraṁ}.) A cup? (\textbf{\ml{കോപ്പ}}, \emph{kōppa}.) A pot? (\textbf{\ml{കലം}}, \emph{kalaṁ}.) A stove? (\textbf{\ml{അടുപ്പ്}}, \emph{aṭuppŭ}.) A hand fan? (\textbf{\ml{വിശറി}}, \emph{viśari}.) An electric fan? (\textbf{\ml{ഫാൻ}}, \emph{phān}.) A bicycle? (\textbf{\ml{സൈക്കിൾ}}, \emph{saikkiḷ}.) A boat? (\textbf{\ml{തോണി}}, \emph{tōṇi}.) A train? (\textbf{\ml{ട്രെയിൻ}}, \emph{ṭreyin}.) A pillow? (\textbf{\ml{തലയണ}}, \emph{talayaṇa}.) A blanket? (\textbf{\ml{പുതപ്പ്}}, \emph{putappŭ}.) A bucket? (\textbf{\ml{ബക്കറ്റ്}}, \emph{bakkaṟṟŭ}.) Scissors? (\textbf{\ml{കത്രിക}}, \emph{katrika}.) A hammer? (\textbf{\ml{ചുറ്റിക}}, \emph{cuṟṟika}.) A nail? (\textbf{\ml{ആണി}}, \emph{āṇi}.) A fence? (\textbf{\ml{വേലി}}, \emph{vēli}.) A gate? (\textbf{\ml{ഗേറ്റ്}}, \emph{gēṟṟŭ}.)
\end{tcolorbox}
